\chapter{MATLAB Code Listings}
\label{app:matlab}

This appendix provides MATLAB implementations for the key algorithms
presented in each chapter. All scripts are self-contained and use only
core MATLAB (no additional toolboxes required).

% ---------------------------------------------------------------
\section{1D Bar Element Assembly and Solution}
\label{lst:bar1d}

\begin{lstlisting}[caption={Three-element bar under tip load (Chapter~\ref{ch:1d})},
                   label={code:bar1d}]
% fem_bar_1d.m -- 3-element 1D bar, tip load
% Example from Chapter 2

clear; clc;

%% Parameters
nElem = 3;              % number of elements
nNode = nElem + 1;      % number of nodes
L     = 0.9;            % total length [m]
EA    = 70e6;           % axial stiffness [N]
F_tip = 10000;          % tip force [N]
Le    = L / nElem;      % element length [m]

%% Element stiffness matrix
ke = (EA/Le) * [1 -1; -1 1];

%% Global assembly
K = zeros(nNode, nNode);
for e = 1:nElem
    dofs = [e, e+1];            % global DOFs for element e
    K(dofs, dofs) = K(dofs, dofs) + ke;
end

%% Force vector
f = zeros(nNode, 1);
f(nNode) = F_tip;               % tip load at last node

%% Apply BC: node 1 fixed (u1 = 0)
freeDoFs = 2:nNode;
K_red = K(freeDoFs, freeDoFs);
f_red = f(freeDoFs);

%% Solve
u_free = K_red \ f_red;
u = [0; u_free];                % full displacement vector

%% Post-processing
fprintf('Nodal displacements [m]:\n');
for i = 1:nNode
    fprintf('  u(%d) = %.6e m\n', i, u(i));
end
\end{lstlisting}

% ---------------------------------------------------------------
\section{Stepped Bar with Distributed and Concentrated Loads}
\label{lst:stepped_bar}

This listing reproduces the worked example of \cref{ex:stepped_bar}: a
two-element cantilever bar with a step in cross-section, a distributed axial
load on the first element, and concentrated forces at nodes~2 and~3.
The script assembles the global stiffness and force vectors element by
element, applies the fixed boundary condition, solves the reduced system,
recovers the reaction, and plots the displacement and internal-force fields
using the enriched trial \cref{eq:bar_disp_enriched}.

\begin{lstlisting}[caption={Stepped bar with distributed and concentrated loads (Chapter~\ref{ch:1d})},
                   label={code:stepped_bar}]
% stepped_bar_2elem.m -- two-element stepped bar (Chapter 3)
% Cantilever bar with two segments:
%   E1: L1 = 500 mm, area  A,  distributed load p = -4 N/mm
%   E2: L2 = 400 mm, area 2A, no distributed load
% Concentrated forces: +6000 N at node 2, -15000 N at node 3.

clear; clc; close all;

%% Material and geometry
E  = 210000;        % Young's modulus [MPa = N/mm^2]
A  = 225;           % element-1 cross-section area [mm^2]
p_e1 = -4;          % distributed load on element 1 [N/mm]
p_e2 =  0;          % distributed load on element 2 [N/mm]

%% Concentrated forces
P_node2 =  6000;    % +x at node 2 [N]
P_node3 = -15000;   % -x at node 3 [N]

%% Discretization
nodes_x   = [0, 500, 900];      % node coordinates [mm]
nElem     = 2;
nNode     = numel(nodes_x);
nDof      = nNode;              % 1D bar: 1 DOF per node
elem_conn = [1 2; 2 3];         % element connectivity
elem_area = [A, 2*A];
elem_p    = [p_e1, p_e2];
elem_L    = diff(nodes_x);

%% Initialize globals
KG    = zeros(nDof, nDof);      % global stiffness
FG_p  = zeros(nDof, 1);         % distributed-load contribution
FG_c  = zeros(nDof, 1);         % concentrated forces

%% Element-by-element assembly
for e = 1:nElem
    L    = elem_L(e);
    Ae   = elem_area(e);
    pe   = elem_p(e);
    dofs = elem_conn(e,:);
    ke   = (E*Ae/L) * [1 -1; -1 1];
    fe_p = (pe*L/2) * [1; 1];
    KG(dofs, dofs) = KG(dofs, dofs) + ke;
    FG_p(dofs)     = FG_p(dofs)     + fe_p;
end

%% Add concentrated forces and form total
FG_c(2) = P_node2;
FG_c(3) = P_node3;
FG = FG_p + FG_c;

%% Boundary conditions: node 1 fixed
fixed_dofs = 1;
free_dofs  = setdiff(1:nDof, fixed_dofs);

UG             = zeros(nDof, 1);
UG(free_dofs)  = KG(free_dofs, free_dofs) \ FG(free_dofs);
RG             = KG*UG - FG;

fprintf('UG = '); disp(UG');
fprintf('RG = '); disp(RG');

%% Element-level postprocessing using the enriched trial:
%   u(x)  = N1*u1 + N2*u2 + p*x*(L-x)/(2EA)
%   N(x)  = EA*(u2-u1)/L + p*(L-2x)/2
xg = []; ug = []; Ng = [];
for e = 1:nElem
    L    = elem_L(e);
    Ae   = elem_area(e);
    pe   = elem_p(e);
    dofs = elem_conn(e,:);
    u1   = UG(dofs(1));  u2 = UG(dofs(2));
    x    = linspace(0, L, 200);
    xi   = x / L;
    u_e  = (1-xi)*u1 + xi*u2 + pe*x.*(L - x)/(2*E*Ae);
    N_e  = E*Ae*(u2-u1)/L + pe*(L - 2*x)/2;
    xg = [xg, nodes_x(dofs(1)) + x];
    ug = [ug, u_e];
    Ng = [Ng, N_e];
end

%% Plots
figure('Color','w'); plot(xg, ug, 'LineWidth', 1.5); grid on;
xlabel('x [mm]'); ylabel('u(x) [mm]');
title('Displacement along the stepped bar');

figure('Color','w'); plot(xg, Ng, 'LineWidth', 1.5); grid on;
xlabel('x [mm]'); ylabel('N(x) [N]');
title('Internal axial force along the stepped bar');
\end{lstlisting}

% ---------------------------------------------------------------
\section{Verification of the Principle of Virtual Work}
\label{lst:pvw_verification}

This listing supports \cref{ex:pvw_verification}: it takes the solved
stepped bar (the same $u_1, u_2, u_3$) and, for three different
kinematically admissible virtual displacements, evaluates the internal
and external virtual work separately. All three cases satisfy
$\delta W_{\mathrm{int}} = \delta W_{\mathrm{ext}}$ to numerical
precision, confirming the principle of virtual work.

\begin{lstlisting}[caption={PVW verification for the stepped bar (Chapter~\ref{ch:1d})},
                   label={code:pvw_verification}]
% pvw_verification_stepped_bar.m
% Numerically verify the principle of virtual work for the SOLVED
% stepped bar (the worked example of Section "Stepped Bar").
% For three different kinematically admissible virtual displacements,
% the script computes the internal and external virtual work separately
% and shows they are equal.

clear; clc;

%% Stepped bar parameters
E  = 210000;
A  = 225;
EA1 = E*A;        % element 1 axial stiffness [N]
EA2 = E*2*A;      % element 2 axial stiffness [N]
L1  = 500;        % element 1 length [mm]
L2  = 400;        % element 2 length [mm]
p1  = -4;         % element 1 distributed load [N/mm]
F2  =  6000;      % concentrated force at node 2 [N]
F3  = -15000;     % concentrated force at node 3 [N]

%% Exact (FEM) nodal displacements from the Stepped Bar example
u1 = 0;            % fixed
u2 = -20/189;      % mm   (= -0.10582)
u3 = -32/189;      % mm   (= -0.16931)

%% Internal axial force in each element using the enriched trial:
%   N(x) = EA*(u_b - u_a)/L  +  p*(L - 2x_local)/2
N1 = @(x) EA1*(u2-u1)/L1 + p1*(L1 - 2*x)/2;        % x measured from node 1
N2 = @(x) EA2*(u3-u2)/L2 + 0*x;                    % x measured from node 2

%% Three admissible virtual displacements (each must satisfy du(x=0)=0)
%   Each entry stores du(x), du'(x), and the nodal values du(node_k)
testCases = struct( ...
  'name',     {}, ...
  'du',       {}, ...
  'du_dx',    {}, ...
  'du_n2',    {}, ...
  'du_n3',    {});

% Case A -- tent: du_2 = 1, du_3 = 0
testCases(end+1) = struct( ...
  'name',  'Case A:  (du_1, du_2, du_3) = (0, 1, 0)', ...
  'du',    @(x_local, elem) (elem==1).*(x_local/L1) + (elem==2).*(1 - x_local/L2), ...
  'du_dx', @(x_local, elem) (elem==1)*(1/L1)         + (elem==2)*(-1/L2), ...
  'du_n2', 1, ...
  'du_n3', 0);

% Case B -- ramp + plateau: du_2 = 1, du_3 = 1
testCases(end+1) = struct( ...
  'name',  'Case B:  (du_1, du_2, du_3) = (0, 1, 1)', ...
  'du',    @(x_local, elem) (elem==1).*(x_local/L1) + (elem==2).*1, ...
  'du_dx', @(x_local, elem) (elem==1)*(1/L1)         + (elem==2)*0, ...
  'du_n2', 1, ...
  'du_n3', 1);

% Case C -- global quadratic du(x_glob) = (x_glob / Ltot)^2
Ltot = L1 + L2;
testCases(end+1) = struct( ...
  'name',  'Case C:  du(x) = (x / 900)^2  (smooth, non-piecewise-linear)', ...
  'du',    @(x_local, elem) ((elem==1).*x_local + (elem==2).*(L1 + x_local)).^2 / Ltot^2, ...
  'du_dx', @(x_local, elem) 2*((elem==1).*x_local + (elem==2).*(L1 + x_local)) / Ltot^2, ...
  'du_n2', (L1/Ltot)^2, ...
  'du_n3', 1);

%% Verify PVW for each case
fprintf('Principle of Virtual Work verification -- stepped bar\n');
fprintf('Solved nodal displacements:  u1 = 0,  u2 = %.5f,  u3 = %.5f mm\n\n', u2, u3);
fprintf('%-58s %14s %14s %14s\n', 'Admissible virtual displacement', ...
        'dW_int [Nmm]', 'dW_ext [Nmm]', '|diff|');
fprintf('%s\n', repmat('-', 1, 102));

for k = 1:length(testCases)
    tc = testCases(k);

    % Internal virtual work -- sum element contributions
    dWint_e1 = integral(@(x) N1(x).*tc.du_dx(x,1), 0, L1);
    dWint_e2 = integral(@(x) N2(x).*tc.du_dx(x,2), 0, L2);
    dWint    = dWint_e1 + dWint_e2;

    % External virtual work -- distributed load (element 1 only) and
    % concentrated forces; reaction at node 1 contributes 0 because du(0)=0.
    dWext_d  = integral(@(x) p1.*tc.du(x,1), 0, L1);
    dWext_c  = F2*tc.du_n2 + F3*tc.du_n3;
    dWext    = dWext_d + dWext_c;

    fprintf('%-58s %14.5f %14.5f %14.2e\n', tc.name, dWint, dWext, abs(dWint-dWext));
end

fprintf(['\nFor every kinematically admissible virtual displacement,\n', ...
         'dW_int = dW_ext to numerical precision: PVW holds.\n']);

\end{lstlisting}

% ---------------------------------------------------------------
\section{2D Beam Element}
\label{lst:beam2d}

\begin{lstlisting}[caption={Simply supported beam with central load (Chapter~\ref{ch:beam})},
                   label={code:beam2d}]
% fem_beam_2d.m -- 2-element Euler-Bernoulli beam
% Example from Chapter 3

clear; clc;

%% Parameters
nElem = 2;
nNode = nElem + 1;
L_total = 2.0;          % total beam length [m]
EI      = 1e4;          % bending stiffness [N*m^2]
P       = 1000;         % central point load [N]
Le      = L_total / nElem;

%% Beam element stiffness (4 DOF: v1,th1,v2,th2)
ke = (EI/Le^3) * [12,   6*Le,  -12,   6*Le;
                   6*Le, 4*Le^2, -6*Le, 2*Le^2;
                  -12,  -6*Le,   12,  -6*Le;
                   6*Le, 2*Le^2, -6*Le, 4*Le^2];

%% DOFs: each node has [v, theta]
nDOF = 2 * nNode;
K = zeros(nDOF, nDOF);

for e = 1:nElem
    n1 = e;  n2 = e + 1;
    dofs = [2*n1-1, 2*n1, 2*n2-1, 2*n2];
    K(dofs, dofs) = K(dofs, dofs) + ke;
end

%% Force vector: central load at node 2 (v direction)
f = zeros(nDOF, 1);
f(2*(nElem/2+1)-1) = P;

%% BCs: pin at node 1 (v1=0) and node 3 (v3=0)
fixedDOFs = [1, 2*nNode-1];
freeDOFs  = setdiff(1:nDOF, fixedDOFs);

K_red = K(freeDOFs, freeDOFs);
f_red = f(freeDOFs);

u_free = K_red \ f_red;
u = zeros(nDOF, 1);
u(freeDOFs) = u_free;

v_mid = u(2*(nElem/2+1)-1);
v_analytic = P * L_total^3 / (48*EI);

fprintf('FEM mid-span deflection:      %.6e m\n', v_mid);
fprintf('Analytical (PL^3/48EI):       %.6e m\n', v_analytic);
fprintf('Error: %.4f%%\n', abs(v_mid - v_analytic)/v_analytic * 100);
\end{lstlisting}

% ---------------------------------------------------------------
\section{CST Element for Plane Stress}
\label{lst:cst2d}

\begin{lstlisting}[caption={Single CST element under plane stress (Chapter~\ref{ch:2d_continuum})},
                   label={code:cst2d}]
% fem_cst_2d.m -- Single CST element, plane stress
% Example from Chapter 4

clear; clc;

E  = 200e9;     % Young's modulus [Pa]
nu = 0.3;       % Poisson's ratio
t  = 0.01;      % thickness [m]

D = (E/(1-nu^2)) * [1,  nu, 0;
                    nu,  1, 0;
                     0,  0, (1-nu)/2];

x = [0; 1; 0];
y = [0; 0; 1];

Ae = 0.5 * abs((x(2)-x(1))*(y(3)-y(1)) - (x(3)-x(1))*(y(2)-y(1)));
fprintf('Element area: %.4f m^2\n', Ae);

b = [y(2)-y(3); y(3)-y(1); y(1)-y(2)];
c = [x(3)-x(2); x(1)-x(3); x(2)-x(1)];

B = (1/(2*Ae)) * [b(1),   0, b(2),   0, b(3),   0;
                    0, c(1),   0, c(2),   0, c(3);
                  c(1), b(1), c(2), b(2), c(3), b(3)];

ke = t * Ae * B' * D * B;

sigma_app = 100e6;
F_total   = sigma_app * 1 * t;
f = zeros(6, 1);
f(3) = F_total / 2;
f(5) = F_total / 2;

freeDOFs = 3:6;
u_free = ke(freeDOFs,freeDOFs) \ f(freeDOFs);
u = zeros(6, 1);
u(freeDOFs) = u_free;

eps_el = B * u;
sig_el = D * eps_el;
fprintf('sigma_x = %.4e\nsigma_y = %.4e\ntau_xy  = %.4e\n', ...
        sig_el(1), sig_el(2), sig_el(3));
\end{lstlisting}

% ---------------------------------------------------------------
\section{Newmark-$\beta$ Time Integration for 2-DOF System}
\label{lst:dynamics}

\begin{lstlisting}[caption={Newmark-$\beta$ transient response of the 2-DOF system (Chapter~\ref{ch:dynamics})},
                   label={code:dynamics}]
% fem_newmark_2dof.m -- Newmark-beta time integration
% 2-DOF mass-spring-damper system from Chapter 5

clear; clc;

m1 = 1500;   m2 = 1000;
k1 = 600000; k2 = 300000;
c1 = 1000;   c2 = 1000;

M = [m1, 0; 0, m2];
C = [c1+c2, -c2; -c2, c2];
K = [k1+k2, -k2; -k2, k2];

beta  = 1/4;
gamma = 1/2;

f_exc  = 2.031;
omega  = 2*pi*f_exc;
t_end  = 5;
dt     = 1e-3;
t_vec  = 0:dt:t_end;
nSteps = length(t_vec);

F_amp = [1000; 1200];
F = @(t) F_amp * sin(omega * t);

u    = zeros(2, 1);
v    = zeros(2, 1);
a    = M \ (F(0) - C*v - K*u);

K_eff = K + (gamma/(beta*dt))*C + (1/(beta*dt^2))*M;

U = zeros(2, nSteps);
U(:,1) = u;

for n = 1:nSteps-1
    f_eff = F(t_vec(n+1)) ...
        + M * (1/(beta*dt^2)*u + 1/(beta*dt)*v + (1/(2*beta)-1)*a) ...
        + C * (gamma/(beta*dt)*u + (gamma/beta-1)*v + dt/2*(gamma/beta-2)*a);

    u_new = K_eff \ f_eff;
    a_new = 1/(beta*dt^2)*(u_new - u) - 1/(beta*dt)*v - (1/(2*beta)-1)*a;
    v_new = v + dt*((1-gamma)*a + gamma*a_new);

    U(:, n+1) = u_new;
    u = u_new;  v = v_new;  a = a_new;
end

figure('Name','Newmark Transient Response');
plot(t_vec, U(1,:)*1000, 'b'); hold on;
plot(t_vec, U(2,:)*1000, 'r');
xlabel('Time [s]'); ylabel('Displacement [mm]');
legend('u_1(t)','u_2(t)'); grid on;
\end{lstlisting}
